\documentclass[12pt,a4paper]{article}
\usepackage[utf8]{inputenc}
\usepackage[T1]{fontenc}
\usepackage[spanish,es-tabla]{babel}
\usepackage{lmodern}
\usepackage[margin=2.5cm]{geometry}
\usepackage{graphicx}
\usepackage{float}
\usepackage{booktabs}
\usepackage{tabularx}
\usepackage{array}
\usepackage{enumitem}
\usepackage{xcolor}
\usepackage{tikz}
\usetikzlibrary{positioning,arrows.meta,calc}
\usepackage{hyperref}
\usepackage{setspace}
\usepackage{fancyhdr}
\setlength{\headheight}{15pt}
\usepackage{pdflscape}
\graphicspath{{../diagramas/}}
\hypersetup{colorlinks=true,linkcolor=blue!50!black,urlcolor=blue!50!black}
\makeatletter\g@addto@macro\UrlBreaks{\do\-}\makeatother
\setstretch{1.1}
\setlength{\parindent}{0pt}
\setlength{\parskip}{6pt}
\newcommand{\pendiente}[1]{\textcolor{red}{\textbf{[POR COMPLETAR: #1]}}}
\pagestyle{fancy}\fancyhf{}
\lhead{Banco de Preguntas Saber}
\rhead{Decisiones iniciales y finales}
\cfoot{\thepage}
\begin{document}
\begin{titlepage}
\centering
{\Large Universidad del Cauca\\Facultad de Ingeniería Electrónica y Telecomunicaciones\\Programa de Ingeniería de Sistemas\par}
\vspace{2cm}
{\huge\bfseries Banco de Preguntas Saber\par}
\vspace{0.5cm}
{\Large Decisiones Iniciales frente a Decisiones Finales\par}
\vspace{1cm}
{Equipo ``Sin Nombre''\par}
\vspace{0.4cm}
Juan Manuel Meneses Daza\\Juan Sebastian Muñoz Ruiz\\Juan Sebastian Molano Bolaños\\Santiago Rodríguez Bolaños
\vfill
Ingeniería de Software II -- 2026.2\\Popayán, noviembre de 2026
\end{titlepage}
\tableofcontents\newpage

\section{Propósito y fuentes}
Este documento compara las decisiones del primer corte con la arquitectura final implementada. La fuente inicial es el PDF de decisiones técnicas del primer corte; las finales se contrastan con los archivos de decisiones de iteración 2 y SAD y el código. La decisión inicial no se interpreta como implementación vigente.

\section{Decisiones iniciales frente a finales}
\begin{landscape}
\small
\renewcommand{\arraystretch}{0.88}
\begin{tabularx}{\linewidth}{p{3.0cm}p{4.4cm}p{5.0cm}X}
\toprule Aspecto & Decisión inicial & Decisión final verificada & Qué cambió y por qué \\
\midrule
Estilo arquitectónico & Monolito en capas con presentación JavaFX, servicios, acceso y dominio. & Cuatro servicios independientes y broker; se conserva \texttt{src/} como cliente/monolito. & El backend se separó para distribuir responsabilidades; cliente no migrado por completo. \\
Plataforma & Java SE y JavaFX con SQLite. & Servicios Spring Boot con SQLite por servicio; JavaFX persiste localmente para preguntas. & Cambia plataforma del backend, preservando legado. \\
Validación & Pipeline de filtros para requisitos estructurales. & Pipeline de siete filtros embebido en Banco. & Se mantiene patrón; se retira servicio independiente para evitar llamada remota en regla local. \\
Mensajería & Observer local para notificar vistas; no microservicios. & RabbitMQ topic \texttt{saberpro.events}; consumidores Revisión y Auditoría. & Observer sigue limitado a GUI; eventos realizan integración distribuida. \\
Ciclo de vida & State modelado por enum en primera versión. & State mediante jerarquía de clases, doble despacho y estados compartidos. & Se localiza comportamiento por estado para representar transiciones permitidas. \\
Cuarto servicio & El diseño intermedio consideraba el servicio de validación independiente. & \texttt{revision-service} (8083) ocupa su lugar; la validación se embebe. & El cambio aplica la sugerencia del profesor del servicio REVISOR a cambio de otro; el motivo exacto se resume en decisiones y código. \\
Auditoría & Historial local en Banco. & \texttt{auditoria-service} proyecta eventos y deduplica por \texttt{eventId}. & Historial separado; consistencia eventual y riesgo de pérdida sin Outbox. \\
Token & Autenticación del monolito. & Token HMAC-SHA256 firmado por Auth, validado localmente con secreto compartido. & Evita consulta de autenticación en cada servicio; exige coordinación de secreto/contrato. \\
Datos & SQLite única para el monolito. & SQLite separada: \texttt{auth.db}, \texttt{banco.db}, \texttt{revision.db}, \texttt{auditoria.db}. & Aísla propiedad de datos, pero limita la demostración de múltiples instancias. \\
Cliente & Aplicación local de escritorio. & Login contra Auth; gestión de preguntas continúa local. & Integración parcial; migración de cliente queda pendiente. \\
\bottomrule
\end{tabularx}
\end{landscape}
\clearpage

\section{Qué se mantuvo}
\begin{itemize}[leftmargin=*]
\item Pipeline como patrón de validación estructural, ahora dentro del Banco.
\item JavaFX como interfaz conservada para la entrega previa; la autenticación cambia a Auth remoto.
\item SQLite se mantiene como tecnología, aunque se separan los archivos por servicio.
\item Principios de dependencias por interfaces en los módulos de dominio y acceso.
\end{itemize}

\section{Decisiones abiertas}
\begin{itemize}[leftmargin=*]
\item Validar con el profesor la aceptación de SQLite respecto al requisito de escalabilidad.
\item Acordar Outbox, recuperación de mensajes y cola de fallidos.
\item Acordar migración de las operaciones del cliente JavaFX al Banco y Revisión remotos.
\item Validar lista de patrones que el curso considera «vistos».
\item Confirmar alcance pendiente de HU05 (asignación y formato de evaluación).
\item Definir gestión del secreto para un entorno no local.
\item \pendiente{El equipo debe aportar evidencias de Jira y aprobación del Product Owner.}
\end{itemize}

\end{document}
